\part{Formal Practice, Automation, and Verification}

\chapter{Constructing Derivations, Models, and Countermodels}
\label{ch:formal-practice}

This chapter collects methods that apply across the logical systems already
defined.  A method is not an inference rule.  It guides the search for a
derivation or semantic object; the completed result must still be checked by
the grammar, calculus, and semantics of the active logical system.

\section{A derivation-search method}
\label{sec:formal-practice-01}

\begin{method}[Backward and forward search]
Given premises \(\Gamma\), conclusion \(\varphi\), and calculus
\(\mathcal C\):
\begin{enumerate}
  \item Copy the exact judgment \(\Gamma\derives_{\mathcal C}\varphi\).
  \item Identify the principal symbol of \(\varphi\), its main
        connective in the propositional case.
  \item Work backward through introduction rules whose conclusion has that
        principal symbol.
  \item Work forward through elimination rules whose major premise is already
        available.
  \item Record every assumption to be discharged and the rule that will
        discharge it.
  \item For a quantifier rule, verify the displayed total capture-avoiding
        substitution and, when present, write the eigenvariable condition
        before applying the rule.
  \item When the branches meet, rewrite the result as a derivation tree and
        check every leaf, rule instance, and discharge mark.
\end{enumerate}
\end{method}

\begin{example}[Mixed search]
To derive \((P\to R)\land(Q\to R)\) from \(P\lor Q\) and
\((P\lor Q)\to R\), backward search on the conjunction creates the two targets
\(P\to R\) and \(Q\to R\).  Each target creates an assumption, to be
discharged by
implication introduction.  Forward search applies the second premise to the
first premise and yields \(R\), independently of either assumption.
Because a discharging application must discharge at least one open
occurrence of its assumption, each implication is then concluded through
a conjunction detour: assume \(P\), conjoin it with the derived \(R\) by
\(\land I\), recover \(R\) by \(\land E_2\), and discharge the occurrence
of \(P\) used by \(\land I\); the case for \(Q\) is symmetric.  The
derivation is longer than
necessary, and that fact is itself useful: the stronger intermediate result
\(R\) should be derived first, after which both implications follow by
the detours just described.
\end{example}

